% ECONOMICS_PROBLEM: {"id": "G14-317", "title": "High-frequency trading and market design", "jel_code": "G14", "source_id": "OP-0099"}
% !TeX program = xelatex
\documentclass[11pt,a4paper]{article}
\usepackage[margin=23mm]{geometry}
\usepackage{fontspec}
\setmainfont{Latin Modern Roman}
\usepackage{amsmath,amssymb,mathtools}
\usepackage{xcolor,enumitem,booktabs,longtable,array,xurl,hyperref}
\definecolor{muted}{HTML}{53636C}
\hypersetup{colorlinks=true,urlcolor=blue,linkcolor=blue}
\setlength{\parindent}{0pt}
\setlength{\parskip}{5pt}
\newcommand{\R}{\mathbb R}
\newcommand{\E}{\mathbb E}
\newcommand{\N}{\mathbb N}
\newcommand{\ind}{\mathbf 1}
\newcommand{\Var}{\operatorname{Var}}
\newcommand{\Cov}{\operatorname{Cov}}
\newcommand{\supp}{\operatorname{supp}}
\DeclareMathOperator*{\argmin}{arg\,min}
\DeclareMathOperator*{\argmax}{arg\,max}
\newcommand{\entrymeta}[3]{{\sffamily\small\color{muted}\textbf{\texttt{#1}}\quad|\quad #2\par #3\par}}
\newcommand{\Problemref}[1]{\texttt{#1}}
\newcommand{\JELref}[2]{\texttt{#2} (JEL \texttt{#1})}
\begin{document}
\section*{High-frequency trading and market design}
\textbf{Problem ID:} \texttt{G14-317}\quad \textbf{JEL:} \texttt{G14}\par
% BEGIN ECONOMICS STATEMENT
\label{problem:OP-0099}
\label{atlas:prob:F9}

\noindent\textbf{Atlas ID:} F9; \textbf{original number:} 99; \textbf{field:} Finance. \textbf{Classification:} Editorial research formulation (theoretical/empirical); not a certified unsolved theorem.

\subsection*{Definitions and assumptions}
A market design $d$ specifies order submission, cancellation, price priority, information release, speed rules and clearing intervals. Traders choose feasible messages, inventories and speed investment; heterogeneous investors demand immediacy under specified preferences. Model $\theta$ gives value innovations, information asymmetry, trader costs and the equilibrium selection for each design, permitting entry and strategic adaptation. Let $W(d;\theta)$ be finite expected investor and trader surplus net of real computing and speed costs, with transfers treated consistently. Define price-discovery loss $V(d)=\mathbb E[(p-v)^2]$ against an explicitly defined efficient-value benchmark $v$, where $p$ is the transaction or quote price under a specified sampling rule.

\subsection*{Formal research question}
For a prespecified implementable set $\mathcal D$, characterize
\[
 d^*\in\operatorname*{arg\,max}_{d\in\mathcal D}W(d;\theta)
 \quad\text{subject to}\quad
 V(d)\le\bar V,\quad\mathbb P_\theta(A(d))\le\bar p,
\]
where $A(d)$ is a defined disruption or execution-failure event and $\bar V,\bar p$ are policy tolerances. Identify changes in liquidity, discovery, rents, resource costs and tail fragility under batch auctions or speed rules. Determine whether empirical comparisons identify equilibrium design effects once trading migrates between venues and traders alter entry, routing and strategies.

\subsection*{Interpretation and scope}
Bid-ask spreads, execution probabilities, adverse-selection losses and price accuracy measure different margins. Speed rents are often transfers, while duplicated technological investment can be a resource cost; adding both without accounting can double count. A common value process and comparable investor demand are necessary for informative counterfactuals. Staggered exchange changes require parallel trends and an explicit interference treatment because securities share liquidity providers and routing networks. Randomization of rules identifies effects for the randomized population but may not identify a market-wide equilibrium redesign. Existence of an optimum requires conditions on the design set, equilibrium correspondence and welfare continuity; strategic discontinuities can violate them.

\subsection*{Source audit and status}
[S1] studies algorithmic trading and liquidity, [S2] proves model-specific implications of frequent batch auctions, and [S3] studies a high-frequency market maker. These are benchmark results and setting-specific evidence rather than unsolved statements. The question is an editorial extension to robust design choice and stressed markets; the audit supplies no exhaustive current-open certification.

\subsection*{References}

\begin{enumerate}[label={[S\arabic*]},leftmargin=*]

\item Terrence Hendershott, Charles M. Jones, and Albert J. Menkveld (2011). \emph{Does Algorithmic Trading Improve Liquidity?}. Journal of Finance 66(1), 1–33. \url{https://doi.org/10.1111/j.1540-6261.2010.01624.x}.
\textit{Locator and role:} Publisher or author abstract / bibliographic record; Algorithmic trading and liquidity around automated quote dissemination. Provides background or a benchmark result; does not establish that the editorial question remains unresolved in 2026.

\item Eric Budish, Peter Cramton, and John Shim (2015). \emph{The High-Frequency Trading Arms Race: Frequent Batch Auctions as a Market Design Response}. Quarterly Journal of Economics 130(4), 1547–1621. \url{https://academic.oup.com/qje/article/130/4/1547/1916146}.
\textit{Locator and role:} Publisher or author abstract / bibliographic record; Speed competition and frequent batch auctions in a specified market-design model. Provides background or a benchmark result; does not establish that the editorial question remains unresolved in 2026.

\item Albert J. Menkveld (2013). \emph{High Frequency Trading and the New Market Makers}. Journal of Financial Markets 16(4), 712–740. \url{https://www.sciencedirect.com/science/article/pii/S1386418113000281}.
\textit{Locator and role:} Publisher or author abstract / bibliographic record; High-frequency market-making behavior in a particular trading setting. Provides background or a benchmark result; does not establish that the editorial question remains unresolved in 2026.

\end{enumerate}

\subsection*{Common conventions: Source mathematical conventions}

Unless an entry states otherwise, agent and firm sets are finite. State and action spaces are measurable, strategies and outcome maps are measurable, probability laws are specified, and expectations exist for the targets under discussion. Infinite-horizon discount factors lie in $(0,1)$ and discounted objectives are finite. Norms, tolerances and complexity input sizes must be specified for computational comparisons. Constants or primitives that appear locally are inputs, not empirical facts asserted by this catalogue.

\textbf{Identification.} A statistical model $m\in\mathcal M$ implies an observable law $P_m$ and target $\theta(m)$. At law $P$, its identified set is
\[
\Theta(P;\mathcal M)=\{\theta(m):m\in\mathcal M,\ P_m=P\}.
\]
Point identification holds when this set is a singleton, or equivalently when $P_{m_1}=P_{m_2}$ implies $\theta(m_1)=\theta(m_2)$ for all compatible models. Sharp bounds mean bounds attained, or approached when the boundary is not attained, by compatible models. Writing this set is a definition, not a solution: the substantive task is to establish informative restrictions and derive or estimate the set. An unrestricted-confounding model may give only trivial bounds.

\textbf{Inference.} Identification, estimability and valid uncertainty quantification are separate questions. Fix $\alpha\in(0,1)$. A confidence procedure $C_n$ over a declared law class $\mathcal P$ has asymptotic uniform coverage at level $1-\alpha$ when
\[
\liminf_{n\to\infty}\inf_{P\in\mathcal P}\Pr_P\{\theta(P)\in C_n\}\geq1-\alpha.
\]
For set-identified targets the coverage event must be defined separately, for example coverage of the full identified set. If an entry uses identification-indexed models instead of uniquely indexed laws, the same coverage convention is applied over those models. Pointwise limits do not establish uniform validity, and asymptotic validity does not guarantee finite-sample precision. Sample sizes, dependence structures and weak-identification sequences are part of each application.

\textbf{Counterfactuals.} $Y_i(a)$ is a potential outcome under a specified intervention, including its timing, implementation and versions. It is not $Y_i$ conditional on voluntarily choosing $a$. A full intervention vector permits interference; shorthand scalar treatment notation does not silently impose no interference. Assignment, selection, attrition, anticipation and measurement must be specified. A claim about an instrument's local effect is not automatically a population policy effect.

\textbf{Economic equilibrium.} A counterfactual model includes best responses, constraints, feasibility, market clearing, state transitions and expectations consistent with the stated information. When several equilibria exist, either a declared selector is an input or the target is set-valued. Comparisons across policy regimes must use the stated selection convention. Reduced-form relationships are not presumed invariant after an equilibrium-changing intervention.

\textbf{Optimization and welfare.} Optimization is over the stated feasible set. If attainment is not established, $\sup$ or $\inf$ and $\varepsilon$-optimal choices replace an asserted maximizer or minimizer. For example, with value $V=\sup_{a\in\mathcal A}W(a)<\infty$, an $\varepsilon$-optimal set is $\{a\in\mathcal A:W(a)\geq V-\varepsilon\}$ for $\varepsilon>0$. Benefits, costs, risk and health or environmental outcomes can be aggregated only using declared units and conversion weights. Interpersonal utility weights, the treatment of fiscal transfers and foreign profits, discounting, and the behavioral welfare criterion are explicit inputs. A comparative-static derivative additionally requires existence and appropriate differentiability of the selected counterfactual.

\textbf{Sources.} Local labels [S1], [S2], and so forth restart in every question. Locators identify the relevant abstract, model, section or result at the level actually checked; a bibliographic or abstract-level verification is not represented as inspection of an entire proof. Working papers are distinguished from later publications and changing versions. Source summaries are paraphrased. All URL checks and editorial formulations refer to the audit date stated above.
% END ECONOMICS STATEMENT
\end{document}
